%AttackName: A-Label Modification-7

%Input:
Let \( x \) be the original input data, \( y \) be the true label, and \( f_{\theta} \) be the target model. The Label Modification attack aims to change the label of an input to mislead the model into making incorrect predictions.

%Output:
The output of the Label Modification attack is a modified label \( y^* \) that causes the model to misclassify the input \( x \).

%Formula:
The Label Modification adversarial attack can be formulated as follows:
1. Initialize the input and true label:
   $
   (x, y)
   $
2. Define the target label \( y^* \) that is different from \( y \):
   $
   y^* \neq y.
   $
3. The goal of the attack is to maximize the likelihood of the incorrect label:
   $
   y^* = \arg\max_{y'} P(f_{\theta}(x) = y' \mid x),
   $
   subject to the constraint that the model's prediction must change:
   $
   f_{\theta}(x) \neq y^*.
   $

%Explanation:
The Label Modification adversarial attack focuses on altering the label associated with an input to deceive the target model. By selecting a target label \( y^* \) that differs from the true label \( y \), the attack aims to maximize the probability of the model predicting this incorrect label when presented with the original input \( x \). This approach highlights the vulnerabilities of models to label manipulation, demonstrating how changing labels can lead to misclassifications. The resulting modified label \( y^* \) serves as an adversarial example that challenges the robustness of the model's predictions.